\section{Gán nhãn vai trò (Semantic Role Labeling)}

\subsection{Đặt vấn đề}

NER xác định \textit{thực thể là gì}, nhưng chưa cho biết \textit{thực thể đó đóng vai gì} trong sự kiện pháp lý. SRL giải quyết câu hỏi \textit{``ai làm gì với ai, khi nào, theo điều kiện nào''} -- thông tin thiết yếu để xây dựng hệ thống hỏi đáp tự động về hợp đồng.

Ví dụ với mệnh đề: \textit{``Party A delivers goods to Party B before May 5, 2024.''}
\begin{itemize}
    \item \textbf{Agent} (ARG0): \textit{Party A} -- người thực hiện hành động
    \item \textbf{Predicate} (V): \textit{delivers}
    \item \textbf{Theme} (ARG1): \textit{goods} -- đối tượng của hành động
    \item \textbf{Recipient} (ARG2): \textit{Party B} -- người nhận
    \item \textbf{Time} (ARGM-TMP): \textit{before May 5, 2024}
\end{itemize}

\subsection{Phương pháp}

\subsubsection{Mô hình}

Sử dụng pretrained model \texttt{yeomtong/srl\_bert\_model} từ HuggingFace Hub:
\begin{itemize}
    \item Checkpoint: \texttt{best\_srl\_Sep\_29.ckpt}
    \item Base model: \texttt{bert-base-cased}
    \item Không cần fine-tune thêm; mô hình dự đoán BIO tags cho từng frame động từ.
\end{itemize}

\subsubsection{Ánh xạ vai nghĩa}

\begin{table}[H]
    \centering
    \caption{Ánh xạ từ PropBank arguments sang vai nghĩa ngữ nghĩa}
    \label{tab:role_map}
    \begin{tabular}{|l|l||l|l|}
        \hline
        \textbf{PropBank} & \textbf{Vai nghĩa} & \textbf{PropBank} & \textbf{Vai nghĩa} \\ \hline
        ARG0     & Agent     & ARGM-TMP & Time      \\ \hline
        ARG1     & Theme     & ARGM-LOC & Location  \\ \hline
        ARG2     & Recipient & ARGM-CAU & Cause     \\ \hline
        ARG3     & Beneficiary & ARGM-ADV & Condition \\ \hline
        ARG4     & Attribute & ARGM-NEG & Negation  \\ \hline
        ARG5+    & Extension & ARGM-MOD & Modal     \\ \hline
    \end{tabular}
\end{table}

\subsubsection{Pipeline xử lý}

\begin{enumerate}
    \item Gọi \texttt{prediction\_formatted(clause\_text)} để lấy BIO tags cho từng frame.
    \item Hàm \texttt{\_decode\_bio\_spans()} giải mã BIO sequence thành các span có vai.
    \item Hàm \texttt{\_normalize\_role\_span()} loại bỏ giới từ dư thừa ở đầu span (ví dụ: \textit{``to Party B''} $\rightarrow$ \textit{``Party B''}).
    \item Chọn \textbf{best frame}: frame có nhiều roles nhất được chọn làm đầu ra chính; tất cả frames được lưu trong \texttt{all\_frames}.
\end{enumerate}

\begin{mintbox}{Pipeline xử lý Semantic Role Labeling}
\begin{minted}[linenos=false, fontsize=\footnotesize]{python}
def srl_pipeline(clause_text):
    # 1. Gọi prediction_formatted để lấy BIO tags cho từng frame
    raw_predictions = prediction_formatted(clause_text)
    # 2. Giải mã BIO sequence thành các spans có vai nghĩa
    decoded_spans = _decode_bio_spans(raw_predictions)
    # 3. Loại bỏ giới từ dư thừa và chuẩn hóa nội dung
    final_roles = {
        role: _normalize_role_span(text) 
        for role, text in decoded_spans.items()
    }
    # 4. Chọn best frame làm đầu ra chính
    return select_best_frame(final_roles)
\end{minted}
\end{mintbox}

\subsection{Kết quả}

\subsubsection{Ví dụ đầu ra SRL}

\begin{minted}[
    linenos=false,
    breaklines,
    frame=single,
    label={Ví dụ kết quả SRL (output/srl\_results.json)}
]{json}
{
  "clause": "The Seller shall transfer legal title of the property to the Buyer.",
  "predicate": "transfer",
  "roles": {
    "Agent": "The Seller",
    "Theme": "legal title",
    "Recipient": "the Buyer"
  },
  "method": "hf_bert_srl"
}
\end{minted}

\subsubsection{Bảng ví dụ nhiều mệnh đề}

\begin{table}[H]
    \centering
    \caption{Kết quả SRL trên các mệnh đề hợp đồng mẫu}
    \label{tab:srl_examples}
    \begin{tabularx}{\textwidth}{|X|l|X|}
        \hline
        \textbf{Mệnh đề} & \textbf{Predicate} & \textbf{Roles} \\ \hline
        Party B shall pay the full rental amount before the 5th of each month.
        & pay
        & Agent: Party B; Theme: rental amount; Time: before the 5th of each month \\ \hline
        The Licensor grants the Licensee the right to use the software.
        & grants
        & Agent: Licensor; Recipient: Licensee; Theme: right to use the software \\ \hline
        If payment is delayed, a penalty of 1\% per day shall apply.
        & apply
        & Theme: penalty; Condition: if payment is delayed \\ \hline
        The Lender may declare the loan balance immediately due and payable.
        & declare
        & Agent: Lender; Theme: loan balance; Modal: may \\ \hline
        The Seller shall transfer legal title of the property to the Buyer.
        & transfer
        & Agent: Seller; Theme: legal title; Recipient: Buyer \\ \hline
    \end{tabularx}
\end{table}

\subsubsection{Trực quan hóa}

\begin{figure}[H]
    \centering
    \begin{dependency}[theme = simple]
        \begin{deptext}[column sep=0.3cm]
            The \& Seller \& shall \& transfer \& legal \& title \& of \& the \& property \& to \& the \& Buyer \\
        \end{deptext}
        \deproot{4}{Predicate (V)}
        \depedge{4}{2}{ARG0}
        \depedge{4}{6}{ARG1}
        \depedge{4}{12}{ARG2}
        
        % Nhóm các từ để thể hiện vai nghĩa là một cụm (span)
        \wordgroup{1}{2}{2}{Agent}
        \wordgroup{1}{5}{6}{Theme}
        \wordgroup{1}{12}{12}{Recipient}
    \end{dependency}
    \caption{Trực quan hóa các semantic roles (spans) trên mệnh đề mẫu mới}
    \label{fig:srl_vis}
\end{figure}

\subsection{Nhận xét và hạn chế}

\begin{itemize}
    \item \textbf{Ưu điểm}: Mô hình pretrained BERT-SRL hoạt động tốt trên văn bản tiếng Anh chuẩn; không cần dữ liệu huấn luyện domain pháp lý.
    \item \textbf{Hạn chế chính}:
    \begin{itemize}
        \item Văn bản hợp đồng pháp lý thường có câu rất dài, nhiều mệnh đề lồng nhau -- model có thể bỏ sót hoặc gán sai role.
        \item Các cấu trúc bị động (\textit{shall be paid by}), modal phức hợp (\textit{may not be required to}) đôi khi không được phân tích đúng.
        \item Khi một mệnh đề có nhiều vị từ, việc chọn ``best frame'' có thể bỏ qua thông tin từ các frames phụ.
    \end{itemize}
\end{itemize}

\begin{figure}[H]
    \centering
    \resizebox{\textwidth}{!}{
        \begin{dependency}[theme = simple]
            \begin{deptext}[column sep=0.15cm]
                Any \& dispute \& arising \& out \& of \& this \& Agreement \& shall \& be \& resolved \& by \& arbitration. \\
            \end{deptext}
            \deproot{10}{V}
            % Lỗi: Chỉ lấy "Any dispute" thay vì cả cụm dài
            \depedge[edge style={red, dashed, thick}, label style={fill=red!10, text=red}]{10}{2}{Theme}
            \depedge{10}{12}{Agent}
            
            \wordgroup[group style={fill=red!10, draw=red, dashed}]{1}{1}{2}{Mô hình chỉ nhận diện cụm ngắn}
            \wordgroup[group style={fill=blue!5, draw=blue, inner sep=2pt}]{1}{1}{7}{Đúng: Toàn bộ chủ thể tranh chấp}
        \end{dependency}
    }
    \caption{Mô hình bỏ sót các thành phần bổ ngữ trong một cụm danh từ dài}
    \label{fig:srl_error}
\end{figure}
